\documentclass{article}
\usepackage{amsmath}
\usepackage{amssymb}

\begin{document}

\begin{titlepage}
    \centering
    \vspace*{6cm} % Adjust the vertical spacing as needed
    {\Huge \textbf{Booleans}} \par
    \vspace*{1cm}
    {\Large Runzhao Guo} \par
    \vspace*{1cm}
    {\Large \today} \par
    \vspace*{2cm}
    % Add a logo or image here using \includegraphics
\end{titlepage}

\section{Question 1}
{(1):}\par
{the truth table of XOR operation is:}\par
\[
\begin{array}{c|c|c}
A & B & A \oplus B \\
\hline
0 & 0 & 0 \\
0 & 1 & 1 \\
1 & 1 & 0 \\
1 & 0 & 1 \\
\end{array}
\]
{when function f and g are the same, their results f(x) and g(x) 
for the same input x are the same. Meaning [f(x), g(x)] can only 
be [0, 0] or [1, 1]. According to the truth table given earlier, 
$h(x) = f(x) \oplus g(x)$ can only output 0}\par
{(2):}\par
{ Assume h(x)'s output is 0 for all xs. We can see that: according 
to the truth Table given in (1) that the input can only be [0, 0] and 
[1, 1]. So the output of f(x) and g(x) are the same to all input x. 
Which means that the two functions are the same function}\par
{combine (1) and (2), we can see that h always outputs 0 if and only 
if f and g are the same function}\par

\section{Question 2}
{(1):}\par
{If $(-1)^x = 1$, 
then 
$\sum_1^n x_n = 2a, a \in \mathbb{Z}$. 
Meaning
$XOR(x) = 2a \mod 2 = 0 , where$ \space $a \in \mathbb{Z}$}\par
{(2):}\par
{ if $XOR(x) = 0$, 
then 
$XOR(x) = 2a \mod 2 = 0 , where$ \space $a \in \mathbb{Z}$ 
Meaning 
$\sum_1^n x_n = 2a, a \in \mathbb{Z}$. 
Which leads to 
$(-1)^{(\sum_1^n x_n)} = (-1)^{2a} = 1, a \in \mathbb{Z}$}\par
{By combining (1) and (2), we can see that $(-1)^x = 1$ holds true 
if and only if $XOR(x) = 0$.}\par

\section{Question 3}
{(1):}\par
{If 
$f(x) = 1$ for all $x$
then
$\forall f(x_1), f(x_2): f(x_1) \oplus f(x_2) = 0$
because
$\forall x, f(x) = 1$
and
$1 \oplus 1 = 0$
}\par
{the same can be said with $f(x) = 0$ for all $x$. 
and that conclusivly prof $f$ is constant if 
$\forall f(x_1), f(x_2): f(x_1) \oplus f(x_2) = 0$}\par

{(2):}\par
{If 
$\forall f(x_1), f(x_2): f(x_1) \oplus f(x_2) = 0$
then 
$\forall f(x_1), f(x_2): f(x_1) = f(x_2)$
Meaning 
$f$ is constant.}\par
{combine (1) and (2), we can prof that $f$ is constant if and only If 
$\forall f(x_1), f(x_2): f(x_1) \oplus f(x_2) = 0$}\par


\end{document}